\section{Desenvolvimento do Módulo de Solicitações}
\label{sec:modulo-solicitacoes}

\subsection{Endpoints}
\label{sec:endpoints}

A camada de rotas do módulo de solicitações foi implementada como um 
\textit{Blueprint} do Flask, registrado sob o prefixo \texttt{/solicitations}. 
A escolha do \textit{Blueprint} decorre diretamente da estratégia de 
\textit{Additive Change} descrita na Seção~\ref{sec:codigo-existente}: 
o módulo é incorporado como um componente isolado, sem alterar os 
controladores existentes. As rotas atuam como camada de fronteira — 
recebem a requisição, aplicam o controle de acesso, delegam o 
processamento ao \texttt{SolicitationService} e serializam a resposta —, 
mantendo fora do controlador as regras de negócio, a persistência e 
o tratamento de exceções de domínio. A documentação Swagger foi 
externalizada em arquivos YAML, referenciados pelo decorador 
\texttt{@swag\_from}, em contraste com o padrão observado nos 
controladores legados (Algoritmo~\ref{alg:legado}), nos quais a 
documentação ocupava a maior parte da função.

O Quadro~\ref{qua:endpoints} relaciona os endpoints implementados, 
a função responsável por cada um, a restrição de acesso aplicada e 
os requisitos funcionais atendidos.

\begin{quadro}[htbp]
\centering
\caption[Endpoints do módulo de solicitações e revisão]
{Endpoints implementados nos módulos de solicitações e revisão.}
\label{qua:endpoints}

\scriptsize
\setlength{\tabcolsep}{4pt}
\renewcommand{\arraystretch}{1.3}

\begin{tabularx}{\linewidth}{
    @{}
    >{\raggedright\arraybackslash}p{1.1cm}
    >{\raggedright\arraybackslash}p{3.9cm}
    >{\raggedright\arraybackslash}X
    >{\raggedright\arraybackslash}p{1.8cm}
    @{}
}
\toprule
\textbf{Método} &
\textbf{Rota} &
\textbf{Controle de acesso} &
\textbf{Requisitos} \\
\midrule

POST & \texttt{/} &
Usuário autenticado &
RF01, RF02, RF03 \\
\midrule

GET & \texttt{/} &
Autenticado com perfil de visualização &
RF04 \\
\midrule

GET & \texttt{/me} &
Usuário autenticado (restrito ao próprio usuário) &
RF04 \\
\midrule

GET & \texttt{/\textless id\textgreater} &
Autenticado, com verificação de acesso no serviço &
RF04, RF11 \\
\midrule

POST & \texttt{/review} &
Usuário autenticado com \texttt{SPECIALIST\_REVIEW} &
RF05, RF06, RF07 \\
\midrule

POST & \texttt{/review/edit} &
Usuário autenticado com \texttt{SPECIALIST\_REVIEW} &
RF05, RF06, RF07 \\
\midrule

GET & \texttt{/review/\textless solicitation\_id\textgreater} &
Usuário autenticado com \texttt{SPECIALIST\_REVIEW} &
RF05, RF06, RF07 \\
\midrule

GET & \texttt{/review/user/\textless user\_id\textgreater} &
Usuário autenticado com \texttt{SPECIALIST\_REVIEW} &
RF05, RF06, RF07 \\

\bottomrule
\end{tabularx}

\fonte{Elaboração própria (2026). Obs.: as rotas do módulo de solicitações são prefixadas com \texttt{/solicitations}, enquanto as rotas de revisão utilizam o prefixo \texttt{/solicitation/review}.}
\end{quadro}


A rota de criação (\texttt{POST /solicitations}) concentra o padrão 
adotado nas demais. O controlador limita-se a extrair os dados da 
requisição — aceitando tanto \texttt{multipart/form-data}, quando há 
arquivo anexado, quanto JSON, quando não há — e a delegar o 
processamento ao serviço, conforme ilustrado no Algoritmo~\ref{alg:rota-criacao}.

\begin{algorithm}[htbp]
\caption{Estrutura da rota de criação de solicitação.}
\label{alg:rota-criacao}
\begin{verbatim}
@bp.route("", methods=["POST"])
@auth.login_required
@swag_from("docs/create_solicitation.yml", validation=False)
def save_solicitation():
    form_data = request.form.to_dict() if request.form else None
    json_data = request.get_json(silent=True) or {}
    file = request.files.get("document_file")

    return SolicitationService.create_or_update(
        user_id=g.user.id,
        form_data=form_data,
        file=file,
    )
\end{verbatim}
\end{algorithm}

Quatro características do padrão merecem destaque. A primeira é o 
\textbf{controlador fino}: a função não constrói entidades, não valida 
campos individualmente, não executa consultas e não trata exceções de 
banco — responsabilidades que, no padrão legado (Algoritmo~\ref{alg:legado}), 
estavam concentradas na própria rota. A segunda é o \textbf{controle 
de acesso declarativo}, aplicado pelo decorador 
\texttt{@auth.login\_required}, que aceita um parâmetro opcional de 
perfil; a rota de listagem geral exige perfil de visualização, enquanto 
a rota \texttt{/me} restringe o resultado ao próprio usuário autenticado, 
independentemente do perfil. A terceira é a \textbf{validação na fronteira} 
por schemas Marshmallow: a consulta é interpretada por 
\texttt{SolicitationQuerySchema}, e a resposta é serializada por 
\texttt{SolicitationSchema} ou \texttt{PaginatedResponseSchema}, de modo 
que o formato de entrada e o de saída não dependem da interface que 
consome a API, atendendo ao RNF04. A quarta é a \textbf{externalização 
da documentação}: o decorador \texttt{@swag\_from} referencia um arquivo 
YAML por rota, com \texttt{validation=False} indicando que a validação 
é responsabilidade dos schemas e não do Swagger.

A listagem foi implementada por meio de uma função auxiliar, 
\texttt{\_fetch\_solicitations}, que centraliza a interpretação dos 
parâmetros de paginação, ordenação e filtro (situação, intervalo de 
datas e usuário). O endpoint \texttt{/solicitations/me} reutiliza a 
mesma função forçando o identificador do usuário autenticado, enquanto 
o endpoint \texttt{/solicitations} permite filtro opcional por usuário, 
desde que o solicitante possua o perfil de visualização. Essa 
centralização evita a duplicação da lógica de consulta entre rotas, 
que, no padrão legado, costumava ser reproduzida em cada controlador.

Do ponto de vista dos requisitos, os endpoints concentram os RF01 a RF04 
(criação e consulta) e o RF11 (restrição de operações a usuários 
autorizados). As operações de decisão, aplicação de efeitos, snapshot e 
reversão, correspondentes aos RF05 a RF09, são implementadas por rotas 
específicas descritas na Seção~\ref{sec:modulo-solicitacoes}.

\subsection{Integração com o Banco de Dados}
\label{sec:banco-dados}

A persistência do módulo foi implementada com SQLAlchemy sobre PostgreSQL, 
seguindo o padrão já adotado pelo sistema. O mapeamento foi feito em estilo 
declarativo, com uma tabela principal (\texttt{solicitations}) e duas tabelas 
de referência (\texttt{solicitation\_status} e \texttt{solicitation\_category}), 
que funcionam como \textit{lookup tables} para os estados e os tipos de 
solicitação. A escolha de manter tabelas de referência, em vez de armazenar 
os valores como texto livre, garante integridade referencial no banco e 
permite a evolução do domínio sem migrations estruturais a cada novo valor.
O Quadro~\ref{qua:modelo-solicitacao} descreve as tabelas do módulo, suas 
colunas e restrições.

\begin{quadro}[htbp]
\centering
\caption[Modelo de dados do módulo de solicitações]{Modelo de dados do módulo de solicitações.}
\label{qua:modelo-solicitacao}
\small
\setlength{\tabcolsep}{3pt}
\begin{tabularx}{\linewidth}{@{}p{2.2cm}p{2.3cm}p{3.4cm}>{\raggedright\arraybackslash}X@{}}
\toprule
\textbf{Tabela} & \textbf{Coluna} & \textbf{Tipo / Restrição} & \textbf{Observação} \\
\midrule
\texttt{solicitations} & \texttt{id} & Inteiro, PK & Identificador da solicitação \\
 & \texttt{id\_user} & FK $\rightarrow$ \texttt{users}, obrigatório & Solicitante \\
 & \texttt{id\_status} & FK $\rightarrow$ \texttt{solicitation\_status}, padrão \texttt{CREATED} & Estado atual \\
 & \texttt{id\_category} & FK $\rightarrow$ \texttt{solicitation\_category}, opcional & Tipo de solicitação; nulo em rascunho \\
 & \texttt{specialist\_type\_id} & FK $\rightarrow$ \texttt{users\_specialist\_types}, opcional & Categoria de destino \\
 & \texttt{title} & Texto(128), obrigatório & Título \\
 & \texttt{description} & Texto(640), obrigatório & Justificativa ou descrição \\
 & \texttt{file\_path} & Texto(256), opcional & Caminho do documento anexado \\
 & \texttt{created\_at} & DataHora, padrão \texttt{agora()} & Criação \\
 & \texttt{last\_update} & DataHora, atualizado em alteração & Última modificação \\
\midrule
\texttt{solicitation status} & \texttt{id}, \texttt{name} & PK, Texto(20) único & \texttt{CREATED}, \texttt{PENDING}, \texttt{APPROVED}, \texttt{REJECTED} \\
\texttt{solicitation category} & \texttt{id}, \texttt{name} & PK, Texto(20) único & \texttt{CHANGE\_CATEGORY}, \texttt{RESOURCE}, \texttt{UPDATE\_POINTS} \\
\texttt{solicitation reviews} & \texttt{id\_solicitation} & FK $\rightarrow$ \texttt{solicitations}, cascata & Decisões e pareceres das análises \\
\bottomrule
\end{tabularx}
\fonte{Elaboração própria (2026).}
\end{quadro}

O Algoritmo~\ref{alg:modelo-solicitation} descreve, em pseudocódigo, a 
estrutura da entidade principal, com seus atributos, restrições, índices 
e relacionamentos, e as enumerações que definem os estados e os tipos de 
solicitação.

\begin{algorithm}[htbp]
\caption{Estrutura da entidade \texttt{Solicitation} e enumerações de domínio (pseudocódigo).}
\label{alg:modelo-solicitation}
\begin{verbatim}
ENUMERAÇÃO SolicitationStatus
    CREATED  = 1    // rascunho, em elaboração
    PENDING  = 2    // submetida, aguardando análise
    APPROVED = 3    // deferida
    REJECTED = 4    // indeferida

ENUMERAÇÃO SolicitationCategory
    CHANGE_CATEGORY = 1    // alteração de categoria
    RESOURCE        = 2    // recurso / complemento cadastral
    UPDATE_POINTS   = 3    // atualização de pontuação

ENTIDADE Solicitation
    TABELA: solicitations

    ATRIBUTOS
        id                 : Inteiro, Chave Primária
        id_user            : Inteiro, FK -> User, obrigatório
        id_status          : Inteiro, FK -> SolicitationStatus,
                             padrão CREATED
        id_category        : Inteiro, FK -> SolicitationCategory,
                             opcional (discriminador de tipo)
        specialist_type_id : Inteiro, FK -> UserSpecialistType,
                             opcional (categoria de destino)
        title              : Texto(128), obrigatório
        description        : Texto(640), obrigatório
        file_path          : Texto(256), opcional
        created_at         : DataHora, padrão agora()
        last_update        : DataHora, atualizado em cada alteração

    RELACIONAMENTOS
        user    : User                 (muitos-para-um)
        reviews : Lista<Review>        (um-para-muitos, cascata)

    ÍNDICES
        idx_solicitation_user     -> (id_user)
        idx_solicitation_status   -> (id_status)
        idx_solicitation_created  -> (created_at)
\end{verbatim}
\end{algorithm}

A decisão de modelagem central é a \textbf{abordagem polimórfica em tabela 
única}: os três tipos de solicitação (\texttt{CHANGE\_CATEGORY}, 
\texttt{RESOURCE} e \texttt{UPDATE\_POINTS}) compartilham a mesma tabela 
\texttt{solicitations}, discriminados pelo campo \texttt{id\_category}. 
Essa escolha evita a fragmentação em três tabelas com colunas comuns 
replicadas (título, descrição, autor, estado, timestamps) e concentra a 
lógica de listagem, filtragem e paginação em uma única consulta, 
independentemente do tipo. Campos específicos de um tipo, como 
\texttt{specialist\_type\_id}, são opcionais e populados apenas quando 
aplicáveis: no fluxo de alteração de categoria, o campo indica a categoria 
de destino; nos fluxos de recurso e de atualização de pontuação, permanece 
nulo. A abordagem foi adotada em vez da herança de tabela por subtipo 
(\textit{joined table inheritance}) por não haver, no escopo do módulo, 
atributos exclusivos de um tipo que justificassem tabelas próprias.

O \textbf{estado de rascunho} é consequência direta dessa modelagem. O 
estado \texttt{CREATED} representa uma solicitação em elaboração, que pode 
ser salva e editada antes da submissão. Nesse estado, o campo 
\texttt{id\_category} pode ser nulo, permitindo que o usuário salve a 
solicitação parcialmente preenchida e escolha o tipo mais adiante. A 
submissão para análise corresponde à transição de \texttt{CREATED} para 
\texttt{PENDING}, momento a partir do qual o tipo passa a ser obrigatório 
na prática, ainda que a restrição não esteja expressa no esquema. A 
alternativa — tornar \texttt{id\_category} obrigatório e criar uma tabela 
separada para rascunhos — foi descartada por duplicar a estrutura e a 
lógica de consulta sem ganho proporcional.

A \textbf{dupla representação de estados e categorias} — enumeração em 
Python e tabela de referência no banco — atende a duas necessidades 
distintas. A enumeração serve à aplicação, garantindo, em tempo de 
desenvolvimento, que apenas valores conhecidos sejam usados e permitindo 
comparações explícitas no código. A tabela de referência serve ao banco, 
garantindo integridade referencial e viabilizando junções e filtros 
eficientes nas consultas de listagem. O custo é a disciplina de manter os 
dois em sincronia a cada novo valor; o benefício é que o estado de uma 
solicitação nunca fica armazenado como texto livre, o que facilitaria 
inconsistências e inviabilizaria restrições declarativas.

A \textbf{chave estrangeira para \texttt{users}}, e não para a entidade de 
especialista, reflete no esquema o desafio identificado na 
Seção~\ref{sec:regras-negocio}: a categoria do especialista está em 
\texttt{users}, e não em \texttt{specialists}. A solicitação referencia 
o usuário e, quando aplicável, indica a categoria de destino em 
\texttt{specialist\_type\_id}, sem duplicar a classificação corrente. Isso 
mantém o modelo coerente com o restante do sistema e evita que a 
solicitação armazene estados intermediários que precisariam ser 
sincronizados com a tabela de origem.

A \textbf{indexação} foi definida a partir das consultas efetivamente 
implementadas nas rotas (Seção~\ref{sec:endpoints}): listagem por usuário 
(\texttt{idx\_solicitation\_user}), listagem por situação 
(\texttt{idx\_solicitation\_status}) e ordenação por data de criação 
(\texttt{idx\_solicitation\_created}). Consultas por identificador já são 
otimizadas pela chave primária e não exigem índice adicional. A escolha 
privilegia os caminhos de leitura da interface — listagem do solicitante, 
listagem do avaliador e ordenação cronológica — em vez de uma indexação 
genérica e não fundamentada.

O \textbf{armazenamento do arquivo} foi delegado ao sistema de arquivos: 
o atributo \texttt{file\_path} guarda o caminho do documento anexado, não 
o conteúdo binário. Essa decisão mantém o banco fora do caminho de 
operações de I/O de arquivos, que são executadas fora da transação 
principal, e evita o inchaço da tabela. A referência ao arquivo permanece 
no mesmo registro da solicitação, mantendo a consistência entre o 
documento e os metadados que o descrevem.

Por fim, o \textbf{controle temporal} foi delegado ao banco: 
\texttt{created\_at} recebe o instante corrente na inserção e 
\texttt{last\_update} é atualizado automaticamente a cada alteração, de 
modo que o momento da última modificação é registrado mesmo em operações 
que não passem pela aplicação. Isso atende parcialmente ao RNF02, no 
aspecto temporal da rastreabilidade; a dimensão semântica — quem alterou, 
de qual estado para qual estado, com qual justificativa — é atendida pelo 
mecanismo de snapshot e pela linha do tempo, descritos na 
Seção~\ref{sec:snapshot-auditoria}. A atomicidade das operações (RNF01), 
por sua vez, não é garantida no modelo, mas na camada de serviço, por 
meio da unidade de trabalho do SQLAlchemy (\texttt{session}): o modelo 
define as restrições; a transação define o momento em que elas são 
verificadas em conjunto.

\subsection{Autorização e Controle de Acesso}

A autorização foi implementada considerando os diferentes papéis envolvidos no fluxo, de
modo que apenas usuários com permissão adequada possam aprovar, rejeitar ou alterar as
informações relacionadas à categoria do especialista. Esse controle é essencial para
preservar as regras de negócio e evitar ações indevidas no sistema.

\subsection{O Módulo de Revisão e a Efetivação das Alterações}

A organização do fluxo de solicitações foi concebida para separar estritamente a
intenção (o pedido do usuário) da efetivação (a decisão e o impacto no cadastro). Para
tornar essa arquitetura compreensível e manter a coesão do relatório, esta seção
divide-se em três aspectos: a separação de responsabilidades e controle de acesso, a
modelagem da revisão como evento de domínio e, por fim, a execução das regras de negócio
e a garantia de atomicidade transacional.

\subsubsection{Separação de Responsabilidades e Controle de Acesso}

O módulo de \textit{Solicitation} atua exclusivamente como o registro da intenção,
enquanto o módulo de \textit{Review} (Revisão) assume a responsabilidade pela validação,
pela decisão e pela efetivação do efeito prático sobre o especialista. A criação da
solicitação apenas documenta o pedido; a decisão autorizada e o impacto material no
cadastro ocorrem no controlador \textit{create\_review}.

Esse componente exige a permissão RBAC \texttt{SPECIALIST\_REVIEW}, restringindo a
execução da aprovação ou rejeição a usuários com competência formal para atuar sobre o
perfil do especialista. A estrutura adotada atende ao requisito de controle de acesso do
sistema e evita que alterações de categoria ou de pontuação sejam aplicadas por atores
sem a devida autorização.

\subsubsection{A Revisão como um Evento de Domínio (Domain Event)}
\label{subsubsec:revisao-como-evento-dominio}
Para assegurar que a decisão seja um registro histórico auditável, a entidade
\texttt{SolicitationReview} foi modelada para atuar como um \textbf{Evento de Domínio}
(\textit{Domain Event}), segundo os preceitos do \textit{Domain-Driven Design} (EVANS,
2003). Em vez de ser apenas uma tabela de passagem, cada registro na tabela
\texttt{solicitation\_reviews} representa o fato imutável de que uma decisão ocorreu no
domínio (o evento de aprovação ou rejeição).

A modelagem armazena a chave estrangeira para a solicitação (\texttt{id\_solicitation}),
o avaliador responsável (\texttt{id\_user}), o veredicto (\texttt{verdict}), a
justificativa obrigatória (\texttt{comment}) e os carimbos temporais
(\texttt{created\_at} e \texttt{last\_update}). A inclusão do usuário garante a dimensão
autoral da auditoria, enquanto o campo \texttt{id\_category} mantém a rastreabilidade
semântica. Adicionalmente, foram definidos índices específicos
(\texttt{idx\_review\_solicitation}, \texttt{idx\_review\_user} e
\texttt{idx\_review\_created}) para otimizar a reconstituição da linha do tempo. Ao
tratar a revisão como um Evento de Domínio, o sistema materializa o gatilho que dispara
os efeitos transacionais no especialista, atendendo diretamente ao RNF02
(Rastreabilidade).

\subsubsection{Execução das Regras de Negócio e Atomicidade}

Antes da efetivação da decisão, o processo de revisão executa validações de consistência
e conflito para garantir que a operação seja compatível com o estado atual do
especialista. O sistema impede, por exemplo, a mudança para a mesma categoria em que o
especialista já se encontra e bloqueia aprovações quando o profissional já está em
processo de retificação. Essas verificações funcionam como guardas de negócio: a decisão
só é aceita quando o contexto do especialista atende às condições mínimas de validade.
Em termos de arquitetura, essa etapa funciona como camada de validação antes da
aplicação dos efeitos, garantindo que os requisitos \textbf{RF05}, \textbf{RF06} e
\textbf{RF07} sejam atendidos de maneira consistente.

Quando a revisão é aprovada (\textit{verdict=True}), a lógica de negócio é delegada à
função interna \textit{\_apply\_review\_business\_rules(solicitation)}, que atua como
uma máquina de estados associada ao tipo da solicitação. Para solicitações do tipo
\texttt{CHANGE\_CATEGORY}, o fluxo aplica a mudança de classificação por meio de
\textit{SpecialistService.switch\_category} e, em seguida, dispara o processo de
retificação via \textit{rectify\_specialist\_service}. Já para solicitações do tipo
\texttt{UPDATE\_POINTS}, o sistema não altera a classificação, mas executa apenas a
retificação do perfil, recalculando a pontuação conforme os critérios definidos.

A etapa de retificação constitui o ponto crítico de atomicidade do processo. A função
\textit{rectify\_specialist\_service} recebe o parâmetro \textit{commit=False}, o que
significa que ela calcula a nova pontuação, altera o status do especialista para
\texttt{RECTIFYING}, registra a ação na linha do tempo e gera a notificação para o
usuário, mas não confirma a transação em si. A confirmação é realizada na camada
superior, por meio de \textit{flush} ou do fechamento explícito da transação, garantindo
que todas as alterações sejam persistidas em conjunto. Esse mecanismo atende ao
requisito \textbf{RNF01}, pois assegura a consistência do estado do especialista e a
impossibilidade de persistir apenas parte das alterações em caso de falha operacional.
Além disso, o uso de \textit{commit=False} preserva a auditoria do sistema: a cronologia
de ações, a notificação enviada ao usuário e a alteração de status ficam vinculadas ao
mesmo evento decisório, permitindo que a responsabilidade pela mudança seja rastreada de
forma inequívoca.
